\chapter{总结与展望}

\section{全文总结}

本文面向低空反无人机场景，以单目可见光相机为唯一传感器，围绕小目标检测难、遮挡跟踪易中断、仿真与真实域差异大三个问题，完成了一套从检测、跟踪、定位到显示的完整系统（DRLS）。按章节顺序回看，各部分工作可以这样归纳。

（1）\textbf{绪论（第1章）}。梳理低空小型无人机的威胁特点与反无人机探测手段的现状，比较雷达、射频与视觉三条技术路线的适用条件，说明本文把单目可见光作为切入点的理由，并给出总体技术路线与章节安排。研究边界也在这一章划定：近程、单目标、可见光。

（2）\textbf{相关理论与数据基础（第2章）}。整理针孔成像与单目测距的几何关系、检测与跟踪的基本理论以及本文采用的评测口径；完成 Det-Fly 的格式转换、统计分析与序列级划分，避免相邻帧之间的数据泄漏；自建 AprilTag 数据集，使用 9$\times$6、25\,mm 棋盘格完成相机标定，为定位精度评估提供可信真值；构建 450 帧仿真自动标注数据集，借助护航相机始终注视目标的特性实现零人工标注，并按距离分层覆盖死区两端。

（3）\textbf{小目标检测与仿真域微调（第3章）}。以 YOLOv8n 为基础，加入 stride=4 的 P2 高分辨率检测头与 EMA 高效多尺度注意力，得到 YOLOv8n-P2-EMA。四组消融实验显示，模型以 2.94\,M 参数量取得 mAP50 83.5\%、Recall 79.2\%、AP\_small 9.8\%，相对基线分别提升 2.2、3.3 与 3.7 个百分点，AP\_small 相对提升 60.7\%，推理速度 98.4\,fps。P2 与 EMA 的分工在这一章得到解释：前者提供“看清”小目标的分辨率，后者提供“认准”小目标的判别力，单独使用都达不到联合后的水平。针对模型在仿真中几乎失效的现象，本文用 2228 帧的距离分层统计与真值几何审计，先以俯仰角分布对比证伪视角假设，再把原因锁定在尺度分布偏移上，随后用 450 帧分层数据微调，将 3--10\,m 的检出率从约 10\% 提到 100\%。

（4）\textbf{三态跟踪与单目定位（第4章）}。设计 SEARCH/LOCKED/COAST 三态状态机，以 0.25 的维持门限与 0.55 的初锁门限组成双门限，抑制误锁；以 140\,px 的空间门限做最近邻关联，抑制身份跳变；以 1.0\,s 的卡尔曼预测维持短时遮挡。参数均有物理依据，算法复杂度为 $O(N)$。定位部分由框中心解算方位角与俯仰角，由框高与等效尺寸解算距离，并推导出测距误差随距离平方增长的规律；六维自适应卡尔曼滤波器采用各向异性量测噪声与 Sage-Husa 在线估计。实验上，TCP 遮挡场景的航迹保持率达到 88.9\%--94.5\%，真实场景录像达到 99.3\%；仿真闭环中单帧测距 RMSE 0.26\,m、AKF 稳态 MedAE 0.13\,m、1\,s 预测 ADE 中位 0.26\,m、方位角误差中位 0.05\textdegree；实拍视频中 AKF 相对单帧测距 RMSE 降低 32\%、抖动降低 33\%。

（5）\textbf{仿真闭环、系统平台与实验（第5章）}。改造 PX4+jMAVSim 闭环，新增脚本目标机、护航相机、真值广播、离屏帧流与城市场景，并用 lockstep 锁住仿真时间；基于 PySide6 实现三页签监视大屏，采用双进程架构与 UDP/JSON 通信，完成态势地图、光电画面、仪表告警、配置持久化与 14 用例自动化测试。实验在同一口径下覆盖检测消融、闭环定位、实拍定位与预测、退化鲁棒性与回归测试，同时划定了距离、尺度、遮挡与光照上的能力边界。

\section{主要创新点}

本文的主要创新点可以归纳为以下四点。

（1）\textbf{面向无人机小目标的 P2+EMA 互补结构及其机理解释}。消融实验验证了改进的有效性，“EMA 单独使用为何失效、与 P2 联合后为何有效”这一问题也得到回答，为小目标检测的网络设计提供了可复用的分析框架。

（2）\textbf{基于真值几何审计的检测死区机理分析方法}。面对“模型在仿真中失效”，本文没有停留在“渲染不真实”的定性解释上，而是以俯仰角分布对比证伪视角假设，以距离分层统计锁定尺度因素，再用微调实验闭环验证。这条“现象—假设—证伪—验证”的路径具有方法论价值。

（3）\textbf{零人工标注的仿真域微调数据构建方法}。护航相机始终注视目标，框中心因此被约束在光轴附近，框高可由真值距离反算，450 帧数据无需人工标注即可生成。该方法成本极低、可扩展性强，适用于任何具备真值广播与可控相机的仿真环境。

（4）\textbf{面向反无人机监视的完整工程实现}。检测、三态跟踪、单目定位、地图态势、光电显示、告警与测试被整合为一套可运行的双进程系统，并在统一口径下完成端到端评估。系统不依赖专用硬件，便于在园区、楼宇与临时活动现场部署。

\section{不足与展望}

\subsection{存在的不足}

（1）\textbf{远距离能力有限}。实拍视频在 10--14\,m 的检出率仅 5.6\%，当前模型与 720p 分辨率的组合覆盖不了更远的告警距离。这既是单目视觉的物理限制，也是系统走向实用的主要瓶颈。

（2）\textbf{相机标定不完整}。实拍评估中相机未做严格标定，焦距按视场角估算，等效尺寸借用仿真值，绝对测距误差达到米级。相对比较仍然有效，但绝对精度不足以支撑实际指挥。

（3）\textbf{单目标假设}。系统面向单目标场景，未处理多架无人机同时出现的情况。画面中一旦出现多个目标，状态机可能在不同目标之间跳变。

（4）\textbf{仿真场景简化}。仿真场景的建筑、树木与光照仍与真实世界有差距，仿真结论不能直接外推到外场；遮挡也主要由静态建筑造成，缺少飞鸟、飘动树叶这类动态遮挡。

（5）\textbf{测试集偏差}。退化实验的测试帧由 simft 模型筛选，对 simft 有利，三组模型的比较不完全公平。

\subsection{未来工作}

（1）\textbf{提升远距离能力}。可以从三个方向入手：提高输入分辨率（如 1280 或 1536）并配合切片推理；换用长焦镜头与可变速云台；引入双目视觉或激光测距辅助深度估计。

（2）\textbf{完善相机标定与尺寸自适应}。用棋盘格完成实拍相机标定，替代视场角估算；按目标机型选择等效尺寸，或利用多帧几何一致性在线估计尺寸。

（3）\textbf{扩展多目标跟踪}。在现有状态机之上增加轨迹池与数据关联层，借鉴 SORT/DeepSORT 的思想，实现多目标的同时跟踪与威胁排序。

（4）\textbf{增强域适应能力}。探索域随机化、风格迁移与自监督预训练，减少对目标域标注的依赖；同时构建更真实的仿真场景，缩小 sim-to-real gap。

（5）\textbf{多传感器融合}。将视觉与射频、声学、雷达结合，由射频承担远距离预警，由视觉完成近程确认与定位，形成互补的探测体系。

（6）\textbf{外场试验}。在合规空域开展真实无人机飞行试验，采集多场景、多机型、多光照的数据，检验系统的实际可用性，并依据外场反馈迭代参数与模型。

\section{结语}

反无人机是一项涉及探测、识别、跟踪、定位与处置的系统工程，单目视觉只是其中的一环。本文的工作表明：针对性的模型改进、严谨的机理分析与完整的工程实现，可以让单目视觉系统在近程反无人机场景中承担实际任务；同时也应看到，单目视觉的物理限制决定了它无法独自完成全部探测工作。未来的反无人机系统必然是多传感器、多算法、多层级协同的体系，本文的工作可以作为其中一个可复现、可扩展的视觉子系统。
